\chapter{Conclusion}

\section{Summary of Findings}\label{sec:findings}

This thesis set out to design and evaluate a drift-aware continuous learning MLOps framework for financial transaction fraud detection that identifies when a deployed model is becoming outdated, decides when retraining is necessary, and deploys a better model safely. Replaying three public datasets under realistic label delay led to five main findings:

\begin{enumerate}
  \item Label-free drift indicators did not reflect model performance: they missed a loss of a third of the PR-AUC on IEEE-CIS and signalled strong drift on Sparkov and BAF, where performance did not fall.
  \item Retraining improved performance on every dataset when its training window suited the data, but label delay eroded its benefit: on IEEE-CIS the gain fell from +0.078 PR-AUC with immediate labels to +0.010 with a 60-day delay.
  \item Tuned fairly, drift-triggered retraining was never reliably better than a fixed schedule and was reliably worse in seven of nine comparisons; with delayed labels the schedule matched or beat it also at an equal retraining budget, while with immediate labels a trigger could retrain more economically.
  \item The training window mattered more than the trigger, and automatic window selection matched the best fixed window on two datasets and came within 0.005 PR-AUC on the third.
  \item The framework reproduced the offline results in a full replay, its promotion gate blocked every faulty model it could observe, and its safety net limited the damage of the fault it could not observe to one week.
\end{enumerate}

\section{Contributions to the Field}\label{sec:contributions}

\begin{enumerate}
  \item An evaluation of retraining policies for fraud detection on public data in strict time order, with explicit label delay, paired bootstrap intervals, walk-forward tuning, repeated seeds and three datasets.
  \item Evidence, on natural drift, that label-free drift indicators are an unreliable trigger for retraining fraud models, and a definition of drift-awareness based on delayed-label performance instead.
  \item An equal-budget comparison showing that, under label delay, the advantage of scheduled retraining is not only due to retraining more often.
  \item A drift-aware MLOps framework with delayed-label monitoring, a scheduled policy with a safety net, automatic window selection and a promotion gate that re-scores models on production features, tested by fault injection.
  \item Open, reproducible code with resumable experiments and Kaggle notebooks, so that every result can be regenerated from the public datasets.
\end{enumerate}

\section{Recommendations for Future Work}\label{sec:future}

\begin{enumerate}
  \item Complete the equal-budget comparison on Sparkov, and load test the scoring service (throughput and p95/p99 latency at increasing concurrency), as recommended in supervisor feedback.
  \item Model faster investigator feedback alongside delayed chargebacks, which could make performance-based triggers react sooner.
  \item Repeat the study with other model families, such as neural and online learners, to test whether the findings depend on gradient boosted trees.
  \item Monitor the stability of model explanations under drift \cite{john2025,uddin2026} and drift in several data sources separately \cite{hassan2026}.
  \item Deploy the framework on a live transaction stream with real label feeds, for example behind a streaming platform \cite{kaushik2025}.
\end{enumerate}
